\editionPage[\pdfbookmark[0]{Corrigés des défis}{bm-defi-cor-1-parent}]{defi-cor-1}{D01 · Corrigé}{Une commande au meilleur prix}{\hyperref[defi-ex-1]{Énoncé p.\,\pageref*{defi-ex-1}}}
\begin{blocEdition}{defi-cor-1}
\editionEtape{Traduire les offres}Chez A : $18\times2{,}40+5=48{,}20$ euros. Chez B : trois lots donnent seulement 15 carnets ; il en faut quatre, soit 20 carnets pour 40 euros.\editionEtape{Comparer}La commande chez B coûte $\res{8{,}20}$ euros de moins et laisse deux carnets supplémentaires. Acheter exactement 18 carnets n’était pas une obligation ; c’est ce qui change le choix.
\end{blocEdition}
\editionPage{defi-cor-2}{D02 · Corrigé}{Deux restes et une divisibilité}{\hyperref[defi-ex-2]{Énoncé p.\,\pageref*{defi-ex-2}}}
\begin{blocEdition}{defi-cor-2}
\editionEtape{Réduire les candidats}Les multiples de 9 sont $108,117,126,135,144$.\editionEtape{Tester la seconde condition}Leurs restes par 7 sont respectivement $3,5,0,2,4$. Seul $\res{135}$ convient : $135=7\times19+2$.\editionEtape{Complétude}La première liste contient tous les multiples de 9 de l’intervalle ; chaque candidat a été testé. Il n’y a donc pas d’autre solution.
\end{blocEdition}
\editionPage{defi-cor-3}{D03 · Corrigé}{Un nombre et son opposé}{\hyperref[defi-ex-3]{Énoncé p.\,\pageref*{defi-ex-3}}}
\begin{blocEdition}{defi-cor-3}
\editionEtape{Utiliser le signe}Si $-a<0$, alors $a$ est strictement positif. Sa distance à zéro, sa valeur absolue, est donc $a$ lui-même.\editionEtape{Combiner les contraintes}On obtient $2<a<5$. Les entiers possibles sont $\res{3\text{ et }4}$. Deux décimaux possibles sont $\res{2{,}5\text{ et }4{,}8}$. Les bornes 2 et 5 sont exclues.
\end{blocEdition}
\editionPage{defi-cor-4}{D04 · Corrigé}{Remonter un carnet de bord}{\hyperref[defi-ex-4]{Énoncé p.\,\pageref*{defi-ex-4}}}
\begin{blocEdition}{defi-cor-4}
\editionEtape{Bilan des déplacements}La variation totale vaut $-35+18{,}5-7{,}5=\res{-24\ \mathrm m}$.\editionEtape{Retrouver le départ}L’altitude initiale vaut $-64-(-24)=\res{-40\ \mathrm m}$.\editionEtape{Contrôler}Depuis $-40$, on passe à $-75$, puis $-56{,}5$, puis $-64$ m. La variation et l’altitude finale ne sont pas la même grandeur.
\end{blocEdition}
\editionPage{defi-cor-5}{D05 · Corrigé}{Une fraction entre deux autres}{\hyperref[defi-ex-5]{Énoncé p.\,\pageref*{defi-ex-5}}}
\begin{blocEdition}{defi-cor-5}
\editionEtape{Encadrer le numérateur}$3/5=9/15$ et $4/5=12/15$, donc $9<a<12$. Les solutions sont $\res{a=10\text{ et }11}$.\editionEtape{Étudier les écritures}$10/15=2/3$ donne une division qui se répète ; $11/15=0{,}7333\ldots$ aussi. Aucune des deux n’a d’écriture décimale finie. Une approximation finie ne doit pas être confondue avec une égalité exacte.
\end{blocEdition}
\editionPage{defi-cor-6}{D06 · Corrigé}{Un ruban en trois parts}{\hyperref[defi-ex-6]{Énoncé p.\,\pageref*{defi-ex-6}}}
\begin{blocEdition}{defi-cor-6}
\editionEtape{Une fraction du tout}$1-2/7-1/3=21/21-6/21-7/21=8/21$.\editionEtape{Retrouver le tout}Huit vingt-et-unièmes valent 32 cm, donc un vingt-et-unième vaut 4 cm. Le ruban mesure $\res{84\ \mathrm{cm}}$.\editionEtape{Répartir}A reçoit 24 cm, B 28 cm, C 32 cm. Le contrôle $24+28+32=84$ confirme le partage.
\end{blocEdition}
\editionPage{defi-cor-7}{D07 · Corrigé}{Le cadre autour d’un carré}{\hyperref[defi-ex-7]{Énoncé p.\,\pageref*{defi-ex-7}}}
\begin{blocEdition}{defi-cor-7}
\editionEtape{Par différence}Le côté extérieur est 10 cm ; $10^2-8^2=\res{36\ \mathrm{cm}^2}$.\editionEtape{Par découpage}Quatre bandes donnent $4\times8\times1=32$ cm$^2$ ; les quatre coins ajoutent 4 cm$^2$. Total : 36 cm$^2$.\editionEtape{Éviter un piège}Quatre bandes de 10 cm sur 1 cm compteraient les coins deux fois.\label{fig-defi-cor-7}\figureGeo[corrige]{puiss-15}

\end{blocEdition}
\editionPage{defi-cor-8}{D08 · Corrigé}{Retrouver l’étape d’un motif}{\hyperref[defi-ex-8]{Énoncé p.\,\pageref*{defi-ex-8}}}
\begin{blocEdition}{defi-cor-8}
\editionEtape{Pour 100}Après retrait de 1, il reste 99, soit 33 groupes de trois. L’étape $\res{33}$ convient.\editionEtape{Pour 101}Il reste 100, qui n’est pas divisible par 3 : aucune étape entière ne convient.\editionEtape{Procédure}Pour un total donné, retirer 1 puis tester la divisibilité par 3. Il faut aussi que l’étape obtenue soit au moins 1.\label{fig-defi-cor-8}\figureGeo[corrige]{lit-8}

\end{blocEdition}
\editionPage{defi-cor-9}{D09 · Corrigé}{La somme de cinq consécutifs}{\hyperref[defi-ex-9]{Énoncé p.\,\pageref*{defi-ex-9}}}
\begin{blocEdition}{defi-cor-9}
\editionEtape{Exprimer}La somme vaut $n+(n+1)+(n+2)+(n+3)+(n+4)=5n+10$.\editionEtape{Factoriser}$5n+10=5(n+2)$. Or $n+2$ est le troisième nombre, celui du milieu.\editionEtape{Conclure}La somme est toujours $\res{\text{cinq fois le nombre du milieu}}$. La lettre rend l’argument valable pour tous les entiers, au-delà d’exemples particuliers.
\end{blocEdition}
\editionPage{defi-cor-10}{D10 · Corrigé}{Un partage dessiné}{\hyperref[defi-ex-10]{Énoncé p.\,\pageref*{defi-ex-10}}}
\begin{blocEdition}{defi-cor-10}
\editionEtape{Retrouver les parts}Les parts seules totalisent $45-6=39$ euros, donc chaque part vaut 13 euros.\editionEtape{Première méthode}La nouvelle part vaut 18 euros ; avec le supplément inchangé, $3\times18+6=\res{60}$ euros.\editionEtape{Seconde méthode}On ajoute trois fois 5 euros : $45+15=60$. Les deux raisonnements donnent le même résultat.\label{fig-defi-cor-10}\figureGeo[corrige]{eq-14}

\end{blocEdition}
\editionPage{defi-cor-11}{D11 · Corrigé}{Deux réductions successives}{\hyperref[defi-ex-11]{Énoncé p.\,\pageref*{defi-ex-11}}}
\begin{blocEdition}{defi-cor-11}
\editionEtape{Calculer les étapes}Après 20 \%, le prix vaut 80 euros. Les 10 \% suivants valent 8 euros ; le prix final est $\res{72}$ euros.\editionEtape{Comparer}La réduction totale vaut $100-72=28$ euros sur 100, soit $\res{28\,\%}$. Additionner 20 \% et 10 \% ne convient pas car les références sont différentes.
\end{blocEdition}
\editionPage{defi-cor-12}{D12 · Corrigé}{Un modèle et son domaine}{\hyperref[defi-ex-12]{Énoncé p.\,\pageref*{defi-ex-12}}}
\begin{blocEdition}{defi-cor-12}
\editionEtape{Calculer}À 3 minutes, il reste 15 L ; à 7 minutes, 3 L. Elle est vide à $\res8$ minutes.\editionEtape{Interpréter}Le tableur continuerait de calculer $24-3t$, donnant $-3$ L à 9 minutes. Le calcul est possible mais ne décrit plus le volume réel. Le modèle doit être arrêté à huit minutes.\label{fig-defi-cor-12}\figureGeo[corrige]{rep-15}

\end{blocEdition}
\editionPage{defi-cor-13}{D13 · Corrigé}{Quatre angles et une bissectrice}{\hyperref[defi-ex-13]{Énoncé p.\,\pageref*{defi-ex-13}}}
\begin{blocEdition}{defi-cor-13}
\editionEtape{Reconstituer l’angle}Il vaut $360-110-95-80=75^\circ$.\editionEtape{Partager}La bissectrice donne $\res{37{,}5^\circ}$ de chaque côté. Elle ajoute une demi-droite à l’intérieur de l’angle, sans déplacer ses côtés ni ceux des autres angles.\label{fig-defi-cor-13}\figureGeo[corrige]{ang-13}

\end{blocEdition}
\editionPage{defi-cor-14}{D14 · Corrigé}{Une aire en plusieurs partages}{\hyperref[defi-ex-14]{Énoncé p.\,\pageref*{defi-ex-14}}}
\begin{blocEdition}{defi-cor-14}
\editionEtape{Grande aire}ABC a une aire de 30 cm$^2$.\editionEtape{Première médiane}CM partage cette aire en deux : ACM a une aire de $\res{15\ \mathrm{cm}^2}$.\editionEtape{Deuxième partage}CN est une médiane du triangle ACM. ACN et NCM ont chacun une aire de $\res{7{,}5\ \mathrm{cm}^2}$. Les aires se déduisent des bases égales et de la hauteur commune.\label{fig-defi-cor-14}\figureGeo[corrige]{defi-triangle-partages}

\end{blocEdition}
\editionPage{defi-cor-15}{D15 · Corrigé}{Un triangle et son centre}{\hyperref[defi-ex-15]{Énoncé p.\,\pageref*{defi-ex-15}}}
\begin{blocEdition}{defi-cor-15}
\editionEtape{Identifier les diagonales}Dans le quadrilatère A-B-A′-B′, les diagonales sont [AA′] et [BB′].\editionEtape{Utiliser le demi-tour}O est le milieu des deux segments.\editionEtape{Conclure}Les diagonales ont le même milieu, donc le quadrilatère est un $\res{\text{parallélogramme}}$. La condition de non-alignement exclut un quadrilatère aplati.\label{fig-defi-cor-15}\figureGeo[corrige]{sym-7}

\end{blocEdition}
\editionPage{defi-cor-16}{D16 · Corrigé}{Même base, plusieurs aires}{\hyperref[defi-ex-16]{Énoncé p.\,\pageref*{defi-ex-16}}}
\begin{blocEdition}{defi-cor-16}
\editionEtape{Aire}Elle reste $8\times3=\res{24\ \mathrm{cm}^2}$ : la formule ne dépend pas du décalage horizontal.\editionEtape{Périmètre}Il n’est pas nécessairement conservé, car les côtés obliques peuvent s’allonger. Par exemple, un rectangle de base 8 et de hauteur 3 a des côtés latéraux de 3 ; après décalage non nul, ils dépassent 3.\editionEtape{Conclusion}Conserver une aire ne suffit pas à conserver un périmètre.\label{fig-defi-cor-16}\figureGeo[corrige]{para-4}

\end{blocEdition}
\editionPage{defi-cor-17}{D17 · Corrigé}{Deux cylindres de même capacité}{\hyperref[defi-ex-17]{Énoncé p.\,\pageref*{defi-ex-17}}}
\begin{blocEdition}{defi-cor-17}
\editionEtape{Volume de référence}$V_A=\pi\times2^2\times10=40\pi$ cm$^3$.\editionEtape{Aire de base de B}Elle vaut $16\pi$ cm$^2$, quatre fois $4\pi$.\editionEtape{Hauteur}Pour le même volume, la hauteur est divisée par quatre : $10:4=\res{2{,}5\ \mathrm{cm}}$. Contrôle : $16\pi\times2{,}5=40\pi$.\label{fig-defi-cor-17}\figureGeo[corrige]{sol-7}

\end{blocEdition}
\editionPage{defi-cor-18}{D18 · Corrigé}{Une moyenne à atteindre}{\hyperref[defi-ex-18]{Énoncé p.\,\pageref*{defi-ex-18}}}
\begin{blocEdition}{defi-cor-18}
\editionEtape{Atteindre 13}La somme connue vaut 48. Il faut $65-48=\res{17}$.\editionEtape{Maximum}Avec 20, la somme vaut 68 et la moyenne $68:5=\res{13{,}6}$. Toute note plus petite donnerait une somme et une moyenne plus petites.
\end{blocEdition}
\editionPage{defi-cor-19}{D19 · Corrigé}{Deux conditions pour une urne}{\hyperref[defi-ex-19]{Énoncé p.\,\pageref*{defi-ex-19}}}
\begin{blocEdition}{defi-cor-19}
\editionEtape{Effectifs connus}Un quart de 20 vaut 5 rouges ; deux cinquièmes valent 8 bleues.\editionEtape{Compléter}Il reste $20-5-8=\res7$ vertes.\editionEtape{Vérifier}Les effectifs sont entiers positifs et totalisent 20. Les probabilités sont $5/20=1/4$, $8/20=2/5$, $7/20$, de somme 1.
\end{blocEdition}
\editionPage{defi-cor-20}{D20 · Corrigé}{Un escalier en boucle}{\hyperref[defi-ex-20]{Énoncé p.\,\pageref*{defi-ex-20}}}
\begin{blocEdition}{defi-cor-20}
\editionEtape{Déterminer la répétition}$150:30=5$ et $100:20=5$ : il faut $\res5$ répétitions.\editionEtape{Distance parcourue}Chaque motif a une longueur de $30+20=50$ pas, donc le lutin parcourt $\res{250}$ pas.\editionEtape{Distinguer les grandeurs}La longueur du trajet en escalier n’est pas la distance en ligne droite entre départ et arrivée.\label{fig-defi-cor-20}\figureGeo[corrige]{prog-12}

\end{blocEdition}
